% !TeX program = xelatex
% !TeX encoding = UTF-8
% !BIB program = biber
% !TeX root = main.tex

\documentclass[
    aspectratio=169,
    xcolor={dvipsnames},
    % handout
]{beamer}
\usetheme[
    % background = dark
    % numbering = none,
]{metropolis}

% \setbeamercolor{background canvas}{bg=white}

% macro.tex

\usepackage{graphicx}
\usepackage{amsmath}
\usepackage{tabularx}
\usepackage{array}
\newcolumntype{x}[1]{>{\centering\arraybackslash\hspace{0pt}}p{#1}}
\usepackage{multicol}
% \setlength{\columnseprule}{1pt}
% \def\columnseprulecolor{\color{block body.fg}}
\usepackage{tikz, pgfplots}
\usetikzlibrary{shapes.geometric, shapes, fit, arrows, calc, fit, positioning, chains, arrows.meta, backgrounds, calc}


% subject specific
\usepackage{amsmath, amssymb}
\usepackage{bussproofs}

% drawing
\usepackage{pstricks}
\usepackage{tikz, tkz-orm, pgfplots}
\usetikzlibrary{arrows,shapes,backgrounds,decorations.markings}
\usetikzlibrary{matrix,positioning,decorations.pathreplacing,calc,tikzmark}
\pgfplotsset{width=7.5cm,compat=1.16}
\usepgfplotslibrary{fillbetween}

% listings
\usepackage[final]{listings}
\usepackage{unicode-math}
\usepackage{xunicode}
\usepackage{fontspec}
\usepackage{pifont}
\setmonofont[
  Contextuals={Alternate}
]{Fira Code}

\definecolor{codegreen}{rgb}{0,0.62,0}
\definecolor{codegray}{rgb}{0.5,0.5,0.5}
\definecolor{codepurple}{rgb}{0.58,0,0.82}
\definecolor{backcolour}{rgb}{1,1,1}

\lstdefinestyle{stolenstyle}{
    backgroundcolor=\color{backcolour},
    commentstyle=\color{codegreen},
    keywordstyle=\color{magenta},
    numberstyle=\tiny\color{codegray},
    stringstyle=\color{codepurple},
    basicstyle=\ttfamily\footnotesize,
    breakatwhitespace=false,
    breaklines=true,
    captionpos=b,
    keepspaces=true,
    numbers=left,
    numbersep=5pt,
    showspaces=false,
    showstringspaces=false,
    showtabs=false,
    tabsize=2,
    inputencoding=utf8,
    extendedchars=true,
    % texcl=false,
    mathescape=true,
    % escapechar=\&,
    literate = 
        {<=>} {3}
        {=>}{{=>}} {2}
        {->}{{->}} {2}
        {_}{{\_}} {1}
        {|}{\textbar} {1}
        {|-}{{|-}} {2}
        {==>}{{==>}} {3},
        % {/\\}{{/\\}} {2}
        % {\\/}{{\\/}} {2}
        % {|-}{{|-}} {2},
}

\lstset{style=stolenstyle}


\lstdefinelanguage{lisa}{
    alsoletter={@,=,>},
    basicstyle=\scriptsize\ttfamily,
    keywordstyle = {\color{blue}},
    keywordstyle = [2]{\color{blue}},
    keywordstyle = [3]{\color{PineGreen}},
    keywordstyle = [4]{\color{teal}},
    commentstyle = \color{Gray},
    morekeywords = [2]{abstract, case, class, def, do, Input, Output, then,
        else, extends, false, free, if, implicit, match,
        object, true, val, var, while, sealed, or,
        for, dependent, null, type, with, try, catch, finally,
        import, final, return, new, override, this, trait,
        private, public, protected, package, throw},
    morekeywords = [3]{have, andThen, thenHave, Theorem, by, DEF, The, Lemma, subproof, assume},
    sensitive = true, 
    numbers=left,
    stepnumber=1,
    morecomment = [l]{//},
    morecomment = [s]{/*}{*/},
    morestring = [b]",  
    otherkeywords = {;,<<,>>,++},
}

\lstset{mathescape,escapeinside={!*}{*!},columns=fullflexible,morekeywords={if,then,else,return,match,with,end,let,in, data, type, :=, def}, basicstyle=\scriptsize\ttfamily}

% \lstset{style=lisa}

% links and citations
\usepackage{hyperref, bookmark}
\hypersetup{hidelinks} % hide boxes around links ew
% pls add color
\usepackage{csquotes}
\usepackage{polyglossia}
\setdefaultlanguage{english}
\usepackage[
    style=numeric,
    sorting=none,
    maxnames=99,
    language=english
]{biblatex}
\addbibresource{biblio.bib}

\usepackage{mathpartir}

\newcommand{\qq}{'\hspace{-0.2em}}



\title{Interpolation and Quantifiers in Orthologic}
\author{Sankalp Gambhir\inst{1}\and Simon Guilloud \inst{1}\and Viktor Kun\v cak \inst{1}}
\institute{\inst{1} Laboratory for Automated Reasoning and Analysis, EPFL}
\date{\today}

\begin{document}

\begin{frame}
    % \titlepage

    \vspace{1em}
    {\huge Interpolation and Quantifiers in Ortholattices}

    \vspace{0.2em}
    {\textcolor[HTML]{DA291C}{\rule{\textwidth}{0.4pt}}}

    \begin{center}
        \footnotesize
        \newcommand{\photowidth}{13ex}
        \begin{tabular}{c c c}
            \includegraphics[width = \photowidth]{authors/sankalp.png} &
            \includegraphics[width = \photowidth]{authors/simon2.jpg}   &
            \includegraphics[width = \photowidth]{authors/viktor.png}  \\
            \emph{Sankalp}                                                    &
            Simon                                                      &
            Viktor                                                     \\
            \emph{Gambhir}                                                    &
            Guilloud                                                   &
            Kun\v cak
        \end{tabular}

    \vspace{5ex}
    \includegraphics[height=6ex]{fig/lara.pdf}
    \end{center}

\end{frame}

\begin{frame}
    \frametitle{Propositional Reasoning}

    % So this begins with the fact that we rely on a lot of propositional
    % reasoning, but it's not easy. Validity, notably, is co-NP hard

    \begin{itemize}[<+->]
        \item We rely extensively on propositional reasoning
        \item Validity checking --- coNP-hard
        \item Can we do better? 
        \item Abstractions?
        \begin{itemize}
            \item Need: soundness
            \item Want: predictability, efficiency 
        \end{itemize}
    \end{itemize}

    % this begs the question, can we reason about propositions in sound and
    % predictable ways, that are not necessarily complete?

    % if we want predictability, perhaps we want to look at algebraic characterizations

\end{frame}

\begin{frame}
    \frametitle{Weakening Boolean Algebras}

    % for one, you have intuitionistic logic, but the validity problem ends up
    % being PSPACE complete

    % another such alternative is orthologic, which, whose models are
    % ortholattices, in the way boolean algebras correspond to propositional logic
    % and heyting algebras to intuitionistic logic.

    % notably, validity in orthologic can be checked in quadratic time

    \begin{itemize}[<+->]
        \item Intuitionistic Logic
        \begin{itemize}
            \item Validity --- PSPACE-Complete
            \footnote<2->{\Citeauthor{statman1979intuitionistic}. \emph{\Citetitle{statman1979intuitionistic}}. In: TCS 1979}
        \end{itemize}
        \item Orthologic
        \begin{itemize}
            \item Validity --- \(O(n^2)\) time
            \footnote<4->{\Citeauthor{guilloud2023formula}. \emph{\Citetitle{guilloud2023formula}}. In: CAV 2023}
        \end{itemize}
    \end{itemize}

    \pause 

    \begin{center}
        \begin{tabular}{c c c}
            Propositional Logic & \(\leftrightarrow\) & Boolean Algebras \\
            Intuitionistic Logic & \(\leftrightarrow\) & Heyting Algebras \\
            \textbf{Orthologic} & \(\leftrightarrow\) & \textbf{Ortholattices}
        \end{tabular}
    \end{center}

\end{frame}

\begin{frame}
    \frametitle{Ortholattices}

    % so how do you define ortholattices?

    % well, starting from the familiar laws for boolean algebras
    % commutativity, associativity, double negation, de Morgan's laws, etc
    % the respective dual laws are not listed due to space

    % the <= can be seen as an implication, written this way because it is a
    % lattice, a poset

    % weaken distributivity to a degenerate version, called absorption


    \begin{center}
        \begin{tabular}{l l}
            Commutativity & \(x \lor y = y \lor x\) \\
            Associativity & \(x \lor (y \lor z) = (x \lor y) \lor z\) \\
            Reflexivity & \(x \lor x = x\) \\
            One & \(x \lor 1 = 1\) \\
            Zero & \(x \lor 0 = x\) \\
            Double Negation & \(\neg \neg x = x\) \\
            Excluded Middle & \(x \lor \neg x = 1\) \\
            De Morgan & \(x \lor (y \land z) = (x \lor y) \land (y \lor z)\) \\
            Distributivity \onslide<2->{\textbf{X}} 
                & \(x \lor (z \land y) = (x \lor z) \land (x \lor y)\) \\
            \pause
            \emph{Absorption} & \(x \lor (x \land y) = x\) 
        \end{tabular}
    \end{center}
\end{frame}

\begin{frame}
    \frametitle{Ortholattices}

    \begin{itemize}
        \item incomplete but sound approximation for Boolean algebras
        \item quadratic-time normalization procedure
        \item found use recently in formula caching for verification
        \footnote{\Citeauthor{guilloud2023formula}. \emph{\Citetitle{guilloud2023formula}}. In: CAV 2023}
        \item as well as in interactive theorem proving \footnote{\Citeauthor{guilloud2023lisa}. \emph{\Citetitle{guilloud2023lisa}}. In: ITP 2023}
    \end{itemize}

\end{frame}

\begin{frame}
    \frametitle{Contributions}

    % In this talk, I will begin by presenting the notion of an ortholattice
    % and since we care primarily about proofs on ortholattices, the corresponding proof system, orthologic
    % and why we care about interpolation in orthologic

    % As an approach to interpolation, I will present our proposed semantics and
    % proof system for quantified orthologic
    % I will talk about the failure of quantifier elimination
    % which makes the question of interpolation more interesting

    % Finally, I will present our results on interpolation for orthologic
    % the failure of a notion of interpolation based on refutation proofs
    % and the existence of interpolation based on implications

    \begin{itemize}
        % \item[\(\checkmark\)] Ortholattices
        % \item Orthologic --- Proof System
        % \item Interpolation for Orthologic
        \item Proof System for Quantified Orthologic (QOL)
        \item Failure of Quantifier Elimination
        \item Failure of Refutation-based Interpolation
        \item Proof of Implicational Interpolation
    \end{itemize}

\end{frame}

\begin{frame}
    \frametitle{Moving to a Sequent Calculus}

      % we are interested in interpolation and proofs, so instead of lattice
      % inequalities, we would like to talk about sequents instead

      % we define sequents as a pair of unordered sets of formulas, the left
      % hand side, the assumptions, and the right hand side, the conclusions

      Sequent
      \begin{gather*}
        \gamma_1, \gamma_2, \ldots \vdash \delta_1, \delta_2, \ldots
      \end{gather*}

      \pause
      Proof rules
      \begin{multicols}{2}
        \centering

        \AxiomC{\phantom{\Gamma}}
        \RightLabel{\text { Hypothesis}}
        \UnaryInfC{$\phi \vdash \phi$}
        \DisplayProof

        \columnbreak

        \centering
        \AxiomC{$\Gamma \vdash \phi$}
        \AxiomC{$\Gamma \vdash \psi$}
        \RightLabel{\text{ RightAnd}}
        \BinaryInfC{$\Gamma \vdash \phi \land \psi$}
        \DisplayProof
      \end{multicols}

      % with a choice of proof rules and other restrictions, if any
      % you get a new proof system of choice

\end{frame}

\begin{frame}
    \frametitle{From Ortholattices to Orthologic}

    % notably, if you restrict classical logic to have atmost one formula on the
    % right hand side, you get intuitionistic logic

    Intuitionistic logic --- at most one formula on the right-hand side of the sequent

    \begin{gather*}
        \gamma_1, \gamma_2, \ldots \vdash \psi
    \end{gather*}

    \pause

    % analogously, orthologic, embodying the laws of ortholattices, is
    % characterized by having atmost two formulas in the sequent at any step

    Orthologic --- at most two formulas in the entire sequent

    \begin{gather*}
        \phi \vdash \psi \\
        \phi, \psi \vdash \\
        \vdash \phi, \psi
    \end{gather*}

    \pause

    % as it's more convenient for orthologic, we write sequents as a set of
    % (atmost two) formulas with left and right annotations

    For convenience, written with left/right annotations, e.g. \(\phi^L, \psi^R\)

\end{frame}

\begin{frame}
    \frametitle{Orthologic --- Proof Rules}
    \centering

    % the proof rules themselves remain the same as for classical propositional
    % sequent calculus, only with the added restriction of having atmost two
    % formulas at any step

    \begin{tabular}{c c}
         \\
            \multicolumn{2}{c}{            
                \AxiomC{}
                \RightLabel{\text { Hyp}}
                \UnaryInfC{$\phi^L, \phi^R$}
                \DisplayProof
            }\\[3ex]

            \multicolumn{2}{c}{
                \AxiomC{$\Gamma, \psi^R$}
                \AxiomC{$\psi^L, \Delta$}
                \RightLabel{\text{ Cut}}
                \BinaryInfC{$\Gamma, \Delta$}
                \DisplayProof
            }\\[3ex]

            \multicolumn{2}{c}{
                \AxiomC{$\Gamma$}
                \RightLabel{\text { Weaken}}
                \UnaryInfC{$\Gamma, \Delta$}
                \DisplayProof
            }\\[3ex]

            
            \AxiomC{$\Gamma, \phi^L$}
            \RightLabel{\text { LeftAnd}}
            \UnaryInfC{$\Gamma, (\phi \land \psi)^L$}
            \DisplayProof &
            \AxiomC{$\Gamma, \phi^R$}
            \AxiomC{$\Gamma, \psi^R$}
            \RightLabel{\text{ RightAnd}}
            \BinaryInfC{$\Gamma, (\phi \land \psi)^R$}
            \DisplayProof
            \\[3ex]

            \AxiomC{$\Gamma, \phi^L$}
            \AxiomC{$\Gamma, \psi^L$}
            \RightLabel{\text{ LeftOr}}
            \BinaryInfC{$\Gamma, (\phi \lor \psi)^L$}
            \DisplayProof &
            \AxiomC{$\Gamma, \phi^R$}
            \RightLabel{\text{ RightOr}}
            \UnaryInfC{$\Gamma, (\phi \lor \psi)^R$}
            \DisplayProof
            \\[3ex]

            \AxiomC{$\Gamma, \phi^R$}
            \RightLabel{\text { LeftNot}}
            \UnaryInfC{$\Gamma, (\neg \phi)^L$}
            \DisplayProof &
            \AxiomC{$\Gamma, \phi^L$}
            \RightLabel{\text{ RightNot}}
            \UnaryInfC{$\Gamma, (\neg \phi)^R$}
            \DisplayProof \\[3ex]
        \end{tabular}

    % pause and breathe

\end{frame}

\begin{frame}
    \frametitle{Orthologic --- Proof Search}

    % following the quadratic normalization procedure, proofs in orthologic are
    % also atmost quadratic in the size of the formulas in the sequent

    % and in fact, even in the presence of arbitrary axioms, proofs can be found
    % in cubic time

    \begin{itemize}
        \item Proofs in orthologic are at most quadratic in size of the formula
        \item Proofs can be found in cubic-time in the presence of axioms\(^1\)
    \end{itemize}

    % you can hear about this, alongside several other interesting results in my
    % coauthors' talk at POPL later this week, titled Orthologic with Axioms 

    \vspace{1em}
    \begin{center}
        \(^1\)\textbf{Orthologic with Axioms} \\
        \emph{Simon Guilloud}, Viktor Kun\v cak\\
        POPL 2024, 19 Jan Friday 16:50
    \end{center}

    % so while incomplete for propositional logic, we can produce small proofs
    % for a large class of formulas

\end{frame}

% Draft notes
% interpolation is useful
%% verification
%% focus model-checking methods on relevant facts
%% starting from proofs, rather than using them simply as certificates of correctness, use them as a basis for generalization 

\begin{frame}
    \frametitle{Interpolation}

    % once we have these proofs, we are interested in interpolation

    % can we find interpolating formulas or "local proofs"

    % interpolation is the question of whether, given say two formulas A and B,
    % and a proof that A implies B, can we find a formula I such that A implies
    % I, and I implies B, while I contains only the free variables common to A
    % and B.

    % in essence, finding the relevant facts in the proof which link the
    % hypothesis and the conclusion

    Given that
    \begin{gather*}
        A_{\bar{x}, \bar{y}} \implies B_{\bar{y}, \bar{z}}
    \end{gather*}

    \pause
    find \(I_{\bar{y}}\) such that
    \(A_{\bar{x}, \bar{y}} \implies I_{\bar{y}}\), and
    \(I_{\bar{y}} \implies B_{\bar{y}, \bar{z}}\).

    \pause
    \begin{itemize}
        \item Focus search to relevant facts % less explosion
        \item Better counterexamples
        \item Abstraction generalization
        \footnote<3->{
            \Citeauthor{mcmillan2018interpolation}. \Citetitle{mcmillan2018interpolation}. In: Handbook of Model Checking
        }
    \end{itemize}

\end{frame}

\begin{frame}
    \frametitle{Interpolation in Boolean Algebras}

    % note that if we allow quantifiers, we get some interpolants for free
    % by simply quantifying away the varaibles that are not common to A and B

    % Boolean algebras, in particular, admit quantifier elimination, which means
    % that for these formulas containing quantifiers (which are not in the
    % language of BA that we need), we can compute equivalent quantifier-free
    % formulas, giving us the interpolants we wanted

    \begin{gather*}
        A_{\bar{x}, \bar{y}} \implies B_{\bar{y}, \bar{z}}
    \end{gather*}

    \pause
    With quantifiers, we get some interpolants for free:

    \begin{itemize}
        \item \(I_1 : \exists \bar{x}. A_{\bar{x}, \bar{y}}\)
        \item \(I_2 : \forall \bar{z}. B_{\bar{y}, \bar{z}}\)
    \end{itemize}

    % we can prove with a little effort that these are indeed interpolants

    \pause
    By quantifier elimination for Boolean algebras, the quantifier-free formulas
    \(QE(I_1)\) and \(QE(I_2)\) can be computed. 

    \pause
    Better: we get least and most general interpolants. \footnote{\Citeauthor{dsilva2010interpolant}. \emph{\Citetitle{dsilva2010interpolant}}. In: VMCAI 2010}
    % also allows you to do some symbolic state computation
\end{frame}


% why is adding quantifiers interesting
%% allows encoding more properties
%% we define semantics of quantified ortholattices

\begin{frame}
    \frametitle{Quantified Orthologic}

    % it would be rgeat if we could do the same for orthologic, but we currently
    % do not have quantifiers at all.

    % so, we introduce semantics of quantified orthologic on complete ortholattices
    % as an analogue to QBF for propositional logic or to system F for intuitionistic logic

    We present the semantics of Quantified Orthologic (QOL) on complete ortholattices.

    % we add two operators, the greatest lower bound, corresponding to universal quantification
    % and the least upper bound, corresponding to existential quantification

    \pause

    \begin{itemize}
        \item Greatest lower bound (universal quantification): \(\bigwedge x. \phi\)
        \item Least upper bound (existential quantification): \(\bigvee x. \phi\)
    \end{itemize}

    % the semantics are given on complete ortholattices to ensure that the bounds really do exist

    \pause

    % correspondingly, we add the usual proof rules for quantification on either
    % side of a sequent
    
    \begin{center}
        \begin{tabular}{c c}
             \\
            \AxiomC{$\Gamma, \phi[x:=\gamma]^L$}
            \RightLabel{\text { LeftForall}}
            \UnaryInfC{$\Gamma, (\bigwedge x. \phi)^L$}
            \DisplayProof &
            \AxiomC{$\Gamma, \phi[x:=x']^R$}
            \RightLabel{\text{ \shortstack{RightForall \\ {\scriptsize($x'$ not free in $\Gamma$)}} }}
            \UnaryInfC{$\Gamma, (\bigwedge x. \phi)^R$}
            \DisplayProof \\[3ex]
    
            \AxiomC{$\Gamma, \phi[x:=x']^L$}
            \RightLabel{\text { \shortstack{LeftExists \\ {\scriptsize($x'$ not free in $\Gamma$)}}}}
            \UnaryInfC{$\Gamma, (\bigvee x. \phi)^L$}
            \DisplayProof &
            \AxiomC{$\Gamma, \phi[x:=\gamma]^R$}
            \RightLabel{\text{ RightExists}}
            \UnaryInfC{$\Gamma, (\bigvee x. \phi)^R$}
            \DisplayProof \\[3ex]
     \\
        \end{tabular}
    \end{center}

\end{frame}

\begin{frame}
    \frametitle{Quantified Orthologic}

    % we show that this new proof system is both sound and complete for
    % ortholattices

    \begin{theorem}[Soundness]
        For every sequent \(S\), if \(\vdash S\) then \(\models S\).
    \end{theorem}

    \begin{theorem}[Completeness]
        For every sequent \(S\), if \(\models S\) then \(\vdash S\).
    \end{theorem}

    % i.e. any inequality provable in QOL is true in all complete ortholattices
    % with quantifiers, and vice versa

    \pause

    Inequalities correspond directly to sequents:
    \begin{gather*}
        \gamma \leq \delta \iff \vdash (\gamma^L, \delta^R)
    \end{gather*}

\end{frame}

% draft notes
% first result: OL does not admit quantifier elimination
%% so quantifiers in OL really do give you more expressive power
%% quantifier elimication in the boolean case give you interpolation
%% does not mean no interpolation

\begin{frame}
    \frametitle{Quantifier Elimination}

    % EMPTY SLIDE WHAT DO I ADD

    % what is qel
    Given a formula \(\phi\) with quantifiers, produce an equivalent
    quantifier-free formula \(QE(\phi)\),

    \begin{gather*}
        \phi \iff QE(\phi)
    \end{gather*}

    % Brown explanation
    % https://www.usna.edu/CS/qepcadweb/B/QE.html
    % exists x phi can be phrased as when does there exist an x such that phi holds?

    
    % \begin{center}
    %     \pause
    %     \(\exists x. \phi\): when does there exist an \(x\) such that \(\phi\) holds?
    
    %     % the quantifier eliminated version is the answer to that question
    %     \pause
    %     There exists an \(x\) such that \(\phi\) holds when... \((\exists x. \phi)_{QE}\)
    % \end{center}
    
    % \begin{center}
    %     \pause
    %     \(\exists x \in \mathbb{R}. a \not = 0 \land ax^2 + bx + c = 0\): \\
    %     when does there exist a real root of the quadratic polynomial?
    
    %     \pause
    %     There exists a real root of the quadratic polynomial when... \(a \not = 0 \land b^2 - 4ac \geq 0\)
    % \end{center}

    % \pause
    % \begin{gather*}
    %     \exists x. a \not = 0 \land ax^2 + bx + c = 0 \iff a \not = 0 \land b^2 - 4ac \geq 0
    % \end{gather*}

\end{frame}

\begin{frame}
    \frametitle{Quantifier elimination in Orthologic}

    % and we show, by counterexample, that quantified orthologic does not admit
    % quantifier elimination. Since the formula phi has only y as a free
    % variable, its quantifier elimination must also contain only y as a free
    % variable.

    % but we can show that any such formula cannot be equivalent to phi
    % simultaneously in both these ortholattice models, M2 and M4. the lattice
    % M2 is distributive, and thus a Boolean algebra, while M4 is not

    \begin{theorem}
        Quantified Orthologic does not admit quantifier elimination.
    \end{theorem}

    \pause

    Counterexample: \(\phi = \bigvee x. (\neg x \land (y \lor x))\)

    \pause
    Quantifier elimination must contain only \(y\): \(0, 1, y, \neg y, \ldots\)

    \pause
    
    \begin{center}
        \fbox{\begin{tikzpicture}[
            roundnode/.style={circle, very thick, minimum size=5mm},
            ]
            %Nodes
            \node[roundnode]      (main)                              {};
            \node[roundnode]      (l)         [above=0.5cm of main] {$1$};
            \node[roundnode]      (o)         [below=0.5cm of main] {$0$};
            \node[roundnode]        (na)       [right=0.2cm of main] {$\neg a$};
            \node[roundnode]        (a)       [left=0.2cm of main] {$a$};
    
            \node[draw = none] (m2) [below of = o] {\(M_2\)};
    
            %Lines
            \draw[-] (l) -- (a);
            \draw[-] (l) -- (na);
            \draw[-] (a) -- (o);
            \draw[-] (na) -- (o);
            \end{tikzpicture}}
            \qquad
                    \fbox{\begin{tikzpicture}[
            roundnode/.style={circle, very thick, minimum size=5mm},
            ]
            %Nodes
            \node[roundnode]      (main)                              {};
            \node[roundnode]      (l)         [above=0.5cm of main] {$1$};
            \node[roundnode]      (o)         [below=0.5cm of main] {$0$};
            \node[roundnode]        (na)       [left=0.2cm of main] {$\neg a$};
            \node[roundnode]        (a)       [left=0.5cm of na] {$a$};
            \node[roundnode]        (b)       [right=0.2cm of main] {$b$};
            \node[roundnode]        (nb)       [right=0.5cm of b] {$\neg b$};
    
            \node[draw = none] (m2) [below of = o] {\(M_4\)};
    
            %Lines
            \draw[-] (l) -- (a);
            \draw[-] (l) -- (na);
            \draw[-] (l) -- (b);
            \draw[-] (l) -- (nb);
            \draw[-] (a) -- (o);
            \draw[-] (na) -- (o);
            \draw[-] (b) -- (o);
            \draw[-] (nb) -- (o);
            \end{tikzpicture}}
    \end{center}

    % so, 1, quantifiers in orthologic really do add more expressive power

    % and 2, it seems we cannot get interpolation so easily, 

\end{frame}

% second result: interpolation 1 does not exist
\begin{frame}
    \frametitle{Refutation-based Interpolation}

    % so we consider two different notions of interpolation, with the first
    % based on refutation proofs.

    Given two sequents \(\gamma\) and \(\delta\), and a proof of contradiction assuming
    \(\gamma\) and \(\delta\), find an interpolant sequent \(I\), such that

    \begin{itemize}[<+->]
        \item \(I\) can be deduced from \(\gamma\) alone
        \item \(I\) and \(\delta\) are inconsistent when assumed together
        \item \(FV(I) \subset FV(\gamma) \cap FV(\delta)\)
    \end{itemize}

    % and we show that orthologic does not admit refutation-based interpolation
    % also by counterexample. 

    % there is an OL proof of contradiction assuming \gamma and \delta, which
    % you can find in the paper 
    
    % Same as quantifier elimination, the interpolant for these two formulas
    % must only contain y as a free variable, and we show that no such formula
    % can be an interpolant for this example.

    \pause
    \begin{theorem}
        Orthologic does not admit refutation-based interpolation.
    \end{theorem}

    \pause

    Counterexample: 
    \begin{gather*}
        \gamma: (z \lor \neg y) \land (\neg z \lor \neg y)^R \\
        \delta: (x \land \neg y) \lor (\neg x \land \neg y)^R
    \end{gather*}
    \pause
    No sequent \(I\) containing only \(y\) can be a refutation-based interpolant
    for this example.

\end{frame}

% third result: interpolation 2 does exist, linear in proof size
\begin{frame}
    \frametitle{Implicational Interpolation for Ortholattices}

    % finally, we consider a notion of interpolation that is more natural for
    % lattices, based on implications, or lattice inequalities, i.e.
    % given two formulas \gamma and \delta, such that \gamma < \delta 

    Given two ortholattice formulas \(\gamma\) and \(\delta\), such that
    \(\gamma \leq \delta\) wrt laws of OL, an implicational interpolant is a
    formula \(I\) such that

    \begin{itemize}
        \item \(\gamma \leq I\)
        \item \(I \leq \delta\)
        \item \(FV(I) \subset FV(\gamma) \cap FV(\delta)\)
    \end{itemize}

\end{frame}

\begin{frame}
    \frametitle{Implicational Interpolation for Ortholattices}

    Given a provable sequent \(\Gamma, \Delta\), an implicational interpolant is
    a formula \(I\) such that

    \begin{itemize}
        \item \(\Gamma, I^R\)
        \item \(I^L, \Delta\)
        \item \(FV(I) \subset FV(\Gamma) \cap FV(\Delta)\)
    \end{itemize}

    \pause

    \begin{theorem}
        Orthologic admits implicational interpolation.
    \end{theorem}

    \pause
    \begin{itemize}
        \item linear-time in size of the proof
        \item thus quadratic-time in size of the sequent
    \end{itemize}
    
\end{frame}

\begin{frame}
    \frametitle{Implicational Interpolation for Ortholattices}

    \begin{theorem}
        Orthologic admits implicational interpolation.
    \end{theorem}

    \begin{proof}
        % base case for hypothesis
        Base case:
        \pause
        
        \begin{center}
            \AxiomC{\phantom{\(\Gamma\)}}
            \RightLabel{\text { Hyp}}
            \UnaryInfC{$\phi^L, \phi^R$}
            \DisplayProof
        \end{center}

        \pause
        Interpolant is \(I = \phi\), and we trivially have proofs of
        \(\phi^L, I^R\) and \(I^L, \phi^R\).
    \end{proof}

\end{frame}

\begin{frame}
    \frametitle{Implicational Interpolation for Ortholattices}

    \begin{theorem}
        Orthologic admits implicational interpolation.
    \end{theorem}

    \begin{proof}
        % base case for hypothesis
        Inductive case:
        \pause

        \begin{center}
            \AxiomC{\(\Gamma, \phi^L\)}
            \RightLabel{\text { LeftAnd}}
            \UnaryInfC{$\Gamma, \phi \land \psi^L$}
            \DisplayProof
        \end{center}

        \pause
        Inductive hypothesis: there is an interpolant \(C\) for \(\Gamma, \phi^L\),
        such that there are proofs of
        \begin{gather*}
            \Gamma, C^R \qquad \qquad C^L, \phi^L
        \end{gather*}

        \vspace{-1em}
        \pause
        We have the interpolant \(I = C\) inductively, and OL proofs of
        interpolation:

        \begin{multicols}{2}
            \centering
            \AxiomC{\(\Gamma, C^R\)}
            \UnaryInfC{\(\Gamma, C^R\)}
            \DisplayProof

            \columnbreak

            \centering
            \AxiomC{\(C^L, \phi^L\)}
            \RightLabel{\text { LeftAnd}}
            \UnaryInfC{$C^L, \phi \land \psi^L$}
            \DisplayProof
        \end{multicols}
    \end{proof}

    % the cases for the other proof rules 

\end{frame}

% possible applications
\begin{frame}
    \frametitle{Orthologic --- Conclusions and Future}

    % starting with
    Starting with
    \begin{itemize}
        \item Orthologic, weakening of classical propositional logic
        \item Sound and complete proof system
    \end{itemize}

    % we show
    We show
    \begin{itemize}
        \item Semantics of Quantified Orthologic
        \item Absence of quantifier elimination
        \item Absence of refutation-based interpolation
        \item Existence of implicational interpolation
        % which is particularly interesting in the absence of qel
    \end{itemize}

    % ortholattices have found use in verification and itp
    % and interpolation has also seen much use in model checking recently
    % we hope these results will lead to more interesting applications of OL in model-checking as well

\end{frame}

\begin{frame}

    \printbibliography    

\end{frame}


\end{document}

% motivation of interpolation

% slow down

% complexity of interpolation


% add future talk to conclusion

% explain ortholattices on the first slide?

% can we get better motivation for ol? something to model with them?

% first slide motivation

% just show the law/rule as different?

% label the law set of boolean algebras

% add citation for intuitionistic logic syntactic restriction

% give example of sequent notations and stop for a bit

% on soundness/completeness side, change provable to words instead















% 1. state contributions

% ^^ 1
% 2. make clear that we will talk about qelim and interpolation

% 3. remove quadratic eqn qel

% 4. make clear that something is an example

% 5. citation for least and greatest

% 6. mcmillan citation, model checking

% 7. swap ol laws and properties

% 8. scale up the laws table, ditch rhs

% 9. remove law labels V*

% 10. quick reminder as to what a sequent is

% 11. shorten quantifier elim def

% 12. does not admit qel, LIKE INTUITIONISTIC LOGIC

% 13. M-2, M-4: M-2 is a BA, M-4 is not

% 14. Give soundness and completeness for QOL formal statements


